% Intended LaTeX compiler: pdflatex
\documentclass[a4paper,12pt]{article}
\usepackage[utf8]{inputenc}
\usepackage[T1]{fontenc}
\usepackage{graphicx}
\usepackage{longtable}
\usepackage{wrapfig}
\usepackage{rotating}
\usepackage[normalem]{ulem}
\usepackage{amsmath}
\usepackage{amssymb}
\usepackage{capt-of}
\usepackage{hyperref}
\usepackage{\string~/.emacs.d/common}
\usepackage{karnaugh-map}
\usepackage{pgfplots}
\usepackage{circuitikz}
\author{mahmood}
\date{\today}
\title{boolean algebra}
\hypersetup{
 pdfauthor={mahmood},
 pdftitle={boolean algebra},
 pdfkeywords={},
 pdfsubject={boolean algebra course},
 pdfcreator={Emacs 28.2 (Org mode 9.5.5)}, 
 pdflang={English}}
\begin{document}

\maketitle
\begin{definition}
a variable can only be \(x \in \{0, 1\}\)\\
\end{definition}
\section{AND}
\label{sec:org9085f40}
\textbf{AND} is a \href{20221219005229-binary_operation.org}{binary operation}\\
\begin{gather*}
  \begin{cases}
    0 \cdot 0 = 0\\
    0 \cdot 1 = 0\\
    1 \cdot 0 = 0\\
    1 \cdot 1 = 1
  \end{cases}
\end{gather*}
\section{OR}
\label{sec:orga2f57d8}
\textbf{OR} is a \href{20221219005229-binary_operation.org}{binary operation}\\
\begin{gather*}
  \begin{cases}
    0 + 0 = 0\\
    0 + 1 = 1\\
    1 + 0 = 1\\
    1 + 1 = 1
  \end{cases}
\end{gather*}
\section{NOT}
\label{sec:orgc9ff093}
\textbf{NOT} is a \href{20221219005258-unary_operation.org}{unary operation}\\
\begin{gather*}
  \begin{cases}
    0' = 1\\
    1' = 0\\
  \end{cases}
\end{gather*}
all of these expressions are equal: \(\lnot x,\overline x,x'\)\\
\section{basic properties of boolean algebra}
\label{sec:orga39f5c1}
\subsection{idempotency}
\label{sec:orgee24117}
\begin{lemma}
\begin{gather*}
  \begin{cases}
    x + x = x\\
    x \cdot x = x
  \end{cases}
\end{gather*}
\end{lemma}
\subsection{identity}
\label{sec:orgf70aa8f}
\begin{lemma}
\begin{gather*}
  \begin{cases}
    A \cdot 1 = A\\
    A + 0 = A
  \end{cases}
\end{gather*}
\end{lemma}
\subsection{annulment}
\label{sec:orgb37efbc}
\begin{lemma}
\begin{gather*}
  \begin{cases}
    A \cdot 0 = 0\\
    A + 1 = 1
  \end{cases}
\end{gather*}
\end{lemma}
\subsection{commutative}
\label{sec:orgd01e956}
\begin{lemma}
\begin{gather*}
  \begin{cases}
    x + y = y + x\\
    x \cdot y = y \cdot x
  \end{cases}
\end{gather*}
\end{lemma}
\subsection{associativity}
\label{sec:org4284a40}
\begin{lemma}
\begin{gather*}
  \begin{cases}
    (x + y) + z = x + (y + z)\\
    (x \cdot y) \cdot z = x \cdot (y \cdot z)\\
  \end{cases}
\end{gather*}
\end{lemma}
\subsection{complement}
\label{sec:org5274b1f}
\begin{lemma}
\begin{gather*}
  \begin{cases}
    x + x' = 1\\
    x \cdot x' = 0
  \end{cases}
\end{gather*}
\end{lemma}
\subsection{distributivity}
\label{sec:orgb7224e0}
\begin{lemma}
\begin{gather*}
  \begin{cases}
    x \cdot (y + z) = x \cdot y + x \cdot z\\
    x + y \cdot z = (x + y) \cdot (x + z)
  \end{cases}
\end{gather*}
\end{lemma}
\subsection{absorption 1}
\label{sec:org3cc7272}
\begin{lemma}
\begin{gather*}
  \begin{cases}
    x + x \cdot y = x\\
    x \cdot (x + y) = x
  \end{cases}
\end{gather*}
\begin{proof}
\begin{align*}
  x + x \cdot y &= x \cdot 1 + x \cdot y\\
                &= x(1 + y)\\
                &= x \cdot 1\\
                &= x
\end{align*}
\begin{align*}
  x \cdot (x + y) &= (x + 0)(x + y)\\
                  &= x + 0 \cdot y\\
                  &= x + 0\\
                  &= x
\end{align*}
\end{proof}
\end{lemma}
\subsection{absorption 2}
\label{sec:orgd64ef6a}
\begin{lemma}
\begin{gather*}
  \begin{cases}
    x + x' \cdot y = x + y\\
    x \cdot (x' + y) = x \cdot y
  \end{cases}
\end{gather*}
\begin{proof}
\begin{align*}
  x + x' \cdot y &= (x + x') \cdot (x + y)\\
                 &= 1 \cdot (x + y)\\
                 &= x + y
\end{align*}
\end{proof}
\end{lemma}
\subsection{consensus}
\label{sec:org283a0d1}
\begin{lemma}
\begin{gather*}
  \begin{cases}
    x \cdot y + x' \cdot z + y \cdot z = x \cdot y + x' \cdot z\\
    (x + y)(x' + z)(y + z) = (x + y)(x' + z)
  \end{cases}
\end{gather*}
\begin{proof}
\begin{align*}
  x \cdot y + x' \cdot z + y \cdot z &= x \cdot y + x' \cdot z + y \cdot z \cdot 1\\
                                     &= x \cdot y + x' \cdot z + y \cdot z \cdot (x + x')\\
                                     &= xy + x'z + yzx + yzx'\\
                                     &= xy(1 + z) + x'z(1 + y)\\
                                     &= xy + x'z
\end{align*}
\end{proof}
\end{lemma}
\subsection{De Morgan}
\label{sec:org6ad57b0}
\begin{lemma}
\begin{note}
not sure yet what the 4th formula refers to\\
\end{note}
\begin{gather*}
  \begin{cases}
    \overline{(\overline{x})} = x\\
    \overline{(x + y)} = \overline{x} \cdot \overline{y}\\
    \overline{(xy)} = \overline{x} + \overline{y}\\
    [T(x_1, x_2, \cdots, x_n, 0, 1, +, \bullet)]' = T(x_1', x_2', \cdots, x_n', 1, 0, \bullet, +)
  \end{cases}
\end{gather*}
\end{lemma}
\section{switching expression}
\label{sec:org177e6eb}
\begin{definition}
a \textbf{switching expression} is a finite combination of switching variables and the operations \hyperref[sec:org9085f40]{AND},\hyperref[sec:orga2f57d8]{OR},\hyperref[sec:orgc9ff093]{NOT}\\
\begin{characteristic}
a combination of switching expressions results in a switching expression, so if \(T_1\) and \(T_2\) are both switching expressions then \(T_1' + T_2\) is also a switching expression\\
\end{characteristic}
\begin{my_example}
an example of simplifying a switching expression\\
\begin{align*}
  T(x, y, z) &= x'y'z + yz + xz\\
             &= z(x'y' + y + x) & \text{distributive property}\\
             &= z(x' + y + x) & \text{absorption 2 property}\\
             &= z(1 + y) & \text{complement property}\\
             &= z \cdot 1 & \text{annulment property}\\
             &= z & \text{identity property}
\end{align*}
\end{my_example}
\begin{my_example}
\begin{align*}
  T(x, y, z) &= (x + y)[x'(y' + z')]' + x'y' + x'z'\\
             &= (x + y)[x + (y' + z')'] + x'y' + x'z'\\
             &= (x + y)[x + yz] + x'y' + x'z'\\
             &= x + yz + x'y' + x'z'\\
             &= x + yz + y' + x'z'\\
             &= x + yz + y' + z'\\
             &= x + z + y' + z'\\
             &= x + 1 + y'\\
             &= 1\\
\end{align*}
\end{my_example}
\end{definition}
\section{switching function}
\label{sec:org4e6eeb6}
a function \(f(x_1, x_2, \cdots, x_n)\) that outputs only \(1\) or \(0\) for any combination of variables it receives\\
\(f\) can only output \(2^n\) different values\\
\begin{my_example}
\begin{tabular}{llll}
\(x_1, \cdots, x_n\) & \(f(x_1, \cdots, x_n)\) & \(f_1\) & \(f_2\)\\
\hline
\(0,\cdots,0\) & \(f(0,\cdots,0)\) & \(1\) & \(0\)\\
\(0, \cdots, 1\) & \(f(0, \cdots, 1)\) & \(0\) & \(0\)\\
\(0, \cdots, 10\) & \(f(0, \cdots, 10)\) & \(0\) & \(1\)\\
\(1,1,1, \cdots, 10\) & \(f(1,1,1, \cdots, 0)\) & \(0\) & \(1\)\\
\(1,1,1, \cdots, 1\) & \(f(1,1,1, \cdots, 1)\) & \(1\) & \(1\)\\
\end{tabular}
\end{my_example}
the number of possible functions for \(n\) variables is \({2^2}^n\)\\
the complement of \(f\), written as \(f'\), is the function that outputs the value 1 wherever \(f\) outputs 0\\
\(f+g\) is the sum of both functions \(g\) and \(f\), \(f + g = 1\) if and only if atleast one of both functions gives \(1\) for a given combination\\
\(f \cdot g\) is the multiplication of both functinos, for a given combination, it gives 1 if and only if both functions give 1 for that specific combination\\
\begin{my_example}
\begin{tabular}{llllllll}
\(x_1\) & \(x_2\) & \(x_3\) & \(f\) & \(f'\) & \(g\) & \(f + g\) & \(f \cdot g\)\\
\hline
\(0\) & \(0\) & \(0\) & \(0\) & \(1\) & \(1\) & \(1\) & \(0\)\\
\(0\) & \(0\) & \(1\) & \(1\) & \(0\) & \(1\) & \(1\) & \(1\)\\
\(0\) & \(1\) & \(0\) & \(1\) & \(0\) & \(0\) & \(1\) & \(0\)\\
\(0\) & \(1\) & \(1\) & \(0\) & \(1\) & \(0\) & \(0\) & \(0\)\\
\(1\) & \(0\) & \(0\) & \(1\) & \(0\) & \(1\) & \(1\) & \(1\)\\
\(1\) & \(0\) & \(1\) & \(0\) & \(1\) & \(0\) & \(0\) & \(0\)\\
\(1\) & \(1\) & \(0\) & \(1\) & \(0\) & \(0\) & \(1\) & \(0\)\\
\(1\) & \(1\) & \(1\) & \(0\) & \(1\) & \(1\) & \(1\) & \(0\)\\
\end{tabular}
\end{my_example}
given \(f\) and \(g\) are switching functions with \(n\) variables, we say \(f\) covers g and write \(g \subseteq f\) if\\
\[
  g(x_1, \cdots, x_n) = 1 \implies f(x_1, \cdots, x_n) = 1
\]\\
if \(g \subseteq f\) and \(g\) is a multiplication of variables then \(g\) is called an of \(f\) and we write \(g \rightarrow f\)\\
\begin{lemma}
\(g \subseteq h + g\)\\
\begin{proof}
\begin{gather*}
  g = 1 \Rightarrow g + h = 1\\
  g \cdot h = 1 \Rightarrow g = 1
\end{gather*}
\end{proof}
\end{lemma}
\begin{lemma}
if \(g \subseteq f\) and \(f \subseteq g\) then \(f = g\)\\
\begin{proof}
\[
  g = 1 \Leftrightarrow f = 1
\]\\
\end{proof}
\end{lemma}
\begin{lemma}
\(g + h \subseteq f\) if and only if \(g \subseteq f\) and \(h \subseteq f\)\\
\begin{proof}
1st side, assume \(g \subseteq f\) and \(h \subseteq f\), we need to prove \(g + h \subseteq f\)\\
assume that \(g + h = 1\), then \(g = 1\) is a consequence of \(g \subseteq f\) that \(f = 1\), and \(h = 1\) is a consequence of \(h \subseteq f\) so that \(f = 1\), and so \(g + h \subseteq f\)\\
2nd side, assume that \(g + h \subseteq f\)\\
if \(g = 1\) then \(g + h = 1\) and so \(f = 1\) and so \(g \subseteq f\)\\
if \(h = 1\) then \(g + h = 1\) and so \(f = 1\) and so \(h \subseteq f\)\\
\end{proof}
\end{lemma}
\subsection{minterm}
\label{sec:org07fe4e3}
\begin{definition}
an expression that is a series of multiplications that contain the variables \(x_1, \cdots, x_n\) of a function \(f(x_1, x_2, \cdots, x_n)\) is called a \textbf{minterm}\\
\end{definition}
\subsection{maxterm}
\label{sec:orgee1eb6c}
\begin{definition}
an expression that is a series of additions that that contain the variables \(x_1, \cdots, x_n\) of a function \(f(x_1, x_2, \cdots, x_n)\) is called a \textbf{maxterm}\\
\end{definition}
\subsection{DNF}
\label{sec:org03ef0dd}
\begin{definition}
when a function is written as a sum of products we get what is called as \textbf{Disjunctive Normal Form}\\
\end{definition}
\subsection{CNF}
\label{sec:org51d7cda}
\begin{definition}
when a function is written as a product of sums we get what is called as \textbf{Conjunctive Normal Form}\\
\end{definition}
\subsection{finding the switching expressions of a switching function}
\label{sec:orgcfcd2f8}
we write a truth table for the function and for each row we write a minterm or a maxterm depending on whether the function outputs 1 or 0, respectively\\
\begin{tabular}{rrrrrll}
 & x & y & z & f & \hyperref[sec:org07fe4e3]{minterm}s & \hyperref[sec:orgee1eb6c]{maxterm}s\\
\hline
0 & 0 & 0 & 0 & 1 & \(x'y'z'\) & \\
1 & 0 & 0 & 1 & 0 &  & \(x+y+z'\)\\
2 & 0 & 1 & 0 & 1 & \(x'yz'\) & \\
3 & 0 & 1 & 1 & 1 & \(x'yz\) & \\
4 & 1 & 0 & 0 & 0 &  & \(x'+y+z\)\\
5 & 1 & 0 & 1 & 0 &  & \(x'+y+z'\)\\
6 & 1 & 1 & 0 & 1 & \(xyz'\) & \\
7 & 1 & 1 & 1 & 1 & \(xyz\) & \\
\end{tabular}
\(f(x, y, z) = \sum(0,2,3,6,7) = x'y'z' + x'yz' + x'yz' + xyz' + xyz \longrightarrow\) sum of products, \hyperref[sec:org03ef0dd]{DNF}\\
\(f(x, y, z) = \Pi(1,4,5) = (x+y+z')(x'+y+z)(x'+y+z') \longrightarrow\) product of sums, \hyperref[sec:org51d7cda]{CNF}\\
\subsection{prime implicant}
\label{sec:org2137fb0}
a multiplication expression \(p\) is called a \textbf{Prime implicant} of a \hyperref[sec:org4e6eeb6]{switching function} \(f\) if:\\
\begin{enumerate}
\item \(p \rightarrow f\) (\(f\) covers \(p\) or \(p\) is an implicant of \(f\))\\
\item deletion of every literal in \(p\) results in product that isnt covered by \(f\) meaning \(p_1, NOT \rightarrow f\) (there is atleast one combination where \(p_1 = 0\) but \(f = 0\))\\
\end{enumerate}
in the last example \(h = x'y\) is a prime implicant of \(f\) because \(x'\) and \(y\) are implicants of \(f\), \(k\) isnt a prime implicant of \(f\)\\
the removal of \(z\), \(x'yz\) turns into \(x'y\) which is still an implicant of \(f\)\\
the removal of \(y \rightarrow x'z\) implicant of \(f\)\\
the removal of \(x' \rightarrow yz\) implicant of \(f\)\\
\begin{lemma}
every sum of products that isnt minimizable that equals \(f\) is a sum of prime implicants of \(f\)\\
\begin{proof}
assume in contradiction that there is a sum of products that isnt minimizable that equals \(f\) and contains the product \(\varphi\) that isnt a prime implicant of \(f\)\\
assume \(f = g + \varphi\), in case of the sum where \(\varphi = 1\) and \(f = 1\) and so \(\varphi \rightarrow f\), in case that \(\varphi\) isnt a prime implicant then its possible to remove from \(\varphi\) atleast 1 literal so that we get \(\varphi_0\) so that \(\varphi_0 \rightarrow f\)\\
also \(\varphi \rightarrow \varphi_0\) in case \(\varphi\) is a product of literals that \(\varphi_0\) is a part of\\
\ldots{}..\\
\end{proof}
\end{lemma}
\subsection{essential prime implicant}
\label{sec:org5d4f600}
from previous theorems\\
\begin{enumerate}
\item a minimal switching expression of a function is the sum of its \emph{Prime Implicants}\\
\item \textbf{Prime Implicants} correspond to squares in the Karnaugh map that arent a part of bigger squares\\
\end{enumerate}
\uline{question}: which \emph{Prime Implicants} exist in the minimal switching expression\\
\uline{definition}: a product of literals \(P\) is called an \emph{Essential Prime Implicant} of the function \(f\) if:\\
\begin{enumerate}
\item \(P\) is a \emph{Prime Implicant} of \(f\)\\
\item \(P\) covers atleast a minterm of \(f\) that cant be covered by another \emph{Prime Implicant}\\
\item \uline{note}: a minimal switching expression has to contain all the possible \emph{EPI's} (\emph{Essential Prime Implicants})\\
\end{enumerate}
\subsection{the process of finding a minimal switching expression of a switching function}
\label{sec:org6111edb}
given the \hyperref[sec:org4e6eeb6]{switching function} \(f\)\\
\begin{itemize}
\item find all the \hyperref[sec:org2137fb0]{prime implicant}s of \(f\)\\
\item from the list of prime implicants, find the \hyperref[sec:org5d4f600]{essential prime implicant}s and put them in an expression\\
\item remove from the list of \emph{PI} all the \emph{EPI}'s and all the \emph{PI}'s that are covered by \emph{EPI}'s\\
\item if the set of \emph{EPI}'s covers \(f\), then we're done, otherwise, choose more \emph{PI}'s so that \(f\) is entirely covered by the resulting set of \emph{PI}'s but make sure to pick a minimal number of \emph{PI}'s to do the job, and make sure to pick the \emph{PI}'s with the least number of literals\\
\end{itemize}
\begin{note}
this isnt a definite solution for the problem but its enough for functions with a small number of literals\\
\end{note}
using karnaugh maps we can apply this process\\
each cell maps to a combination of variables, we place each minterm in its corresponding cell, i.e. with the cell that provides the correct corresponding values for the variables\\
\includesvg{/home/mahmooz/brain/out/cqS6PbM}\\

\includesvg{/home/mahmooz/brain/out/nixxAs4}\\

\includesvg{/home/mahmooz/brain/out/WJLkEpu}\\

by finding rectangles whose area is a power of 2 in these tables we would be finding minterms and the least amount of rectangles we can find would give us the minimal form we're looking for\\
the variables that cells in a rectangle may share are the ones to be taken for the switching expression\\

\begin{my_example}
\begin{gather*}
  f(w,x,y,z) = \sum(4,5,8,12,13,14,15)\\
  f = xy' + wx + wy'z'
\end{gather*}
\begin{figure}[htbp]

\includesvg{/home/mahmooz/brain/out/dQOeJio}
\caption{\emph{Prime implicants}: \(w\bar{y}\bar{z}, wx, x\bar{y}\), all of them are \emph{EPI}}
\end{figure}
\end{my_example}
\begin{my_example}
\begin{gather*}
  f(x,y,z) = \sum(0,2,3,4,5,7)
\end{gather*}
\begin{figure}[htbp]

\includesvg{/home/mahmooz/brain/out/96j2WfN}
\caption{\emph{Prime Implicants}: \(\bar{x}\bar{z}, \bar{x}y, yz, xz, x\bar{y}, \bar{y}\bar{z}\)}
\end{figure}

none of these \emph{PI}'s isnt \emph{EPI} and all of them are the same size\\
this function has 2 minimal forms\\
\(f(x,y,z) = \bar{x}y + xz + \bar{y}\bar{z}\)\\
\(f(x,y,z) = \bar{x}\bar{z} + yz + x\bar{y}\)\\
\end{my_example}
\subsection{dont care}
\label{sec:orgd886573}
\textbf{dont care} is a value of a function we get for a combination of values for the variables that isnt defined, and therefore can be either \(1\) or \(0\)\\
the symbol of the value \emph{Don't care} is \(\phi\), its allowed to switch the value \(\phi\) with either \(1\) or \(0\) according to our needs, this is useful when finding the minimal form of a switching function\\
\subsection{map expressions}
\label{sec:org06f3bb0}
\begin{definition}
a map expression is a switching expression that appears in the map and that isnt \(0\), \(1\) or \(\phi\)\\
\begin{characteristic}
the usage of map exxpressions allows us to describe a function with more than \(n\) variables using a karnaugh map with \(n\) variables\\
\end{characteristic}
\begin{my_example}
consider the following function \(f(x,y,z,A,B,C)\)\\
\begin{tabular}{rrrrrrr}
x & y & z & A & B & C & f\\
\hline
0 & 0 & 0 & 0 & 0 & 0 & 0\\
0 & 0 & 0 & 0 & 0 & 1 & 0\\
0 & 0 & 0 & 0 & 1 & 0 & 0\\
0 & 0 & 0 & 0 & 1 & 1 & 0\\
0 & 0 & 0 & 1 & 0 & 0 & 0\\
0 & 0 & 0 & 1 & 0 & 1 & 0\\
0 & 0 & 0 & 1 & 1 & 0 & 0\\
0 & 0 & 0 & 1 & 1 & 1 & 0\\
0 & 0 & 1 & 0 & 0 & 0 & 0\\
0 & 0 & 1 & 0 & 0 & 1 & 0\\
0 & 0 & 1 & 0 & 1 & 0 & 1\\
0 & 0 & 1 & 0 & 1 & 1 & 1\\
0 & 0 & 1 & 1 & 0 & 0 & 0\\
0 & 0 & 1 & 1 & 0 & 1 & 0\\
0 & 0 & 1 & 1 & 1 & 0 & 1\\
0 & 0 & 1 & 1 & 1 & 1 & 1\\
0 & 1 & 0 & 0 & 0 & 0 & 0\\
0 & 1 & 0 & 0 & 0 & 1 & 0\\
0 & 1 & 0 & 0 & 1 & 0 & 0\\
0 & 1 & 0 & 0 & 1 & 1 & 0\\
0 & 1 & 0 & 1 & 0 & 0 & 1\\
0 & 1 & 0 & 1 & 0 & 1 & 1\\
0 & 1 & 0 & 1 & 1 & 0 & 1\\
0 & 1 & 0 & 1 & 1 & 1 & 1\\
0 & 1 & 1 & 0 & 0 & 0 & 1\\
0 & 1 & 1 & 0 & 0 & 1 & 1\\
0 & 1 & 1 & 0 & 1 & 0 & 1\\
0 & 1 & 1 & 0 & 1 & 1 & 1\\
0 & 1 & 1 & 1 & 0 & 0 & 1\\
0 & 1 & 1 & 1 & 0 & 1 & 1\\
0 & 1 & 1 & 1 & 1 & 0 & 1\\
0 & 1 & 1 & 1 & 1 & 1 & 1\\
1 & 0 & 0 & 0 & 0 & 0 & 0\\
1 & 0 & 0 & 0 & 0 & 1 & 1\\
1 & 0 & 0 & 0 & 1 & 0 & 0\\
1 & 0 & 0 & 0 & 1 & 1 & 1\\
1 & 0 & 0 & 1 & 0 & 0 & 0\\
1 & 0 & 0 & 1 & 0 & 1 & 1\\
1 & 0 & 0 & 1 & 1 & 0 & 0\\
1 & 0 & 0 & 1 & 1 & 1 & 1\\
1 & 0 & 1 & 0 & 0 & 0 & 1\\
1 & 0 & 1 & 0 & 0 & 1 & 1\\
1 & 0 & 1 & 0 & 1 & 0 & 0\\
1 & 0 & 1 & 0 & 1 & 1 & 0\\
1 & 0 & 1 & 1 & 0 & 0 & 1\\
1 & 0 & 1 & 1 & 0 & 1 & 1\\
1 & 0 & 1 & 1 & 1 & 0 & 0\\
1 & 0 & 1 & 1 & 1 & 1 & 0\\
1 & 1 & 0 & 0 & 0 & 0 & 0\\
1 & 1 & 0 & 0 & 0 & 1 & 0\\
1 & 1 & 0 & 0 & 1 & 0 & 0\\
1 & 1 & 0 & 0 & 1 & 1 & 0\\
1 & 1 & 0 & 1 & 0 & 0 & 1\\
1 & 1 & 0 & 1 & 0 & 1 & 1\\
1 & 1 & 0 & 1 & 1 & 0 & 1\\
1 & 1 & 0 & 1 & 1 & 1 & 1\\
1 & 1 & 1 & 0 & 0 & 0 & \(\phi\)\\
1 & 1 & 1 & 0 & 0 & 1 & \(\phi\)\\
1 & 1 & 1 & 0 & 1 & 0 & \(\phi\)\\
1 & 1 & 1 & 0 & 1 & 1 & \(\phi\)\\
1 & 1 & 1 & 1 & 0 & 0 & \(\phi\)\\
1 & 1 & 1 & 1 & 0 & 1 & \(\phi\)\\
1 & 1 & 1 & 1 & 1 & 0 & \(\phi\)\\
1 & 1 & 1 & 1 & 1 & 1 & \(\phi\)\\
\end{tabular}
we can represent it using a 3 variable karnaugh map\\
\[ A’BC’ + A’BC + ABC’ + ABC = B(A’C’+A’C+AC’+AC) = B \]\\

\includesvg{/home/mahmooz/brain/out/N9R81Eh}\\
\end{my_example}
\end{definition}
\section{functional completeness}
\label{sec:orgd41edb1}
\begin{definition}
a set of operations forms a \textbf{functionally complete} system if every switching function can be described as a switching expression that only contains operations from that set\\
\begin{my_example}
for example \(\{\lnot, +\}\) is functionally complete because \(x \cdot y = (x'+y')'\) and in every switching expression where we have \(x \cdot y\) we can replace it with \(\lnot(\lnot x + \lnot y)\)\\
\end{my_example}
\begin{my_example}
\(\{\cdot, \lnot\}\) is also functionally complete\\
\end{my_example}
\begin{lemma}
a functionally complete system consisting of only 1 function returns the complement of a variable if its passed to all its inputs, and the opposite is true, as in a system that doesnt abide by this rule isnt functionally complete\\
\begin{my_example}
given the completely functional system \(\{f(x,y,z,w)\}\), according to this rule, \(f(x,x,x,x) = \overline{x}\)\\
\end{my_example}
\begin{my_example}
take the functionally complete operation \emph{NOR}, we pass \(x\) to all its inputs and we get \(\overline{x}\):\\
\[ x \downarrow x = \overline{x} \cdot \overline{x} = \overline{x} \]\\
\end{my_example}
\end{lemma}
\begin{lemma}
one way to prove that a set of operations is completely functional is by describing all the operations of a set that is already known to be functionally complete using the operations in that set\\
\begin{my_example}
proof of NOR's functional completeness using the already functionally complete set \(\{\lnot, +\}\):\\
\begin{gather*}
  x \downarrow x = x' \cdot x' = x'\\
  (x \downarrow y) \downarrow (x \downarrow y) = (x' \cdot y')' \cdot (x' \cdot y')' = (x + y)(x + y) = (x + y)\\
  ((x + y)' + (x + y)')' = (x'y' + x'y')' = (x'y')' \cdot (x'y')' = (x+y)(x+y) = (x+y)
\end{gather*}
\end{my_example}
\end{lemma}
\end{definition}
\subsection{list of systems known to be functionally complete}
\label{sec:org8ef39a2}
\begin{itemize}
\item \(\{+, \cdot, \lnot\}\)\\
\item \(\{+, \lnot\}\)\\
\item \(\{\cdot, \lnot\}\), because \(x + y = \overline{(\bar{x} \cdot \bar{y})}\)\\
\item \emph{NOR}: \(x \downarrow y = \overline{x} \cdot \overline{y} = \overline{(x + y)}\)\\
\item \emph{NAND}: \(x \uparrow y = \overline{x} + \overline{y} = \overline{(xy)}\)\\
\end{itemize}
\subsection{partial functional completeness}
\label{sec:org62859a2}
\begin{definition}
a system of operations is \textbf{partially functionally complete} if it need to be combined with \(0\) and \(1\) or both to provide functional completeness\\
\begin{note}
even if a set of operations can produce 1's and 0's, its still gonna be partially complete if we substitute 1 or 0 to provide completeness\\
\begin{my_example}
\[
f_1(x,y) = x \oplus y \qquad f_2(x,y) = x + y'
\]\\
\(f_1(x,y)=x'y+xy'\)\\
\(f_1(x,1)=x' \longrightarrow\) derived negation\\
\(f_2(x, f_2(0, y)) = x+y \longrightarrow\) derived OR operator\\
so this system is partially complete\\
\end{my_example}
\end{note}
\end{definition}
\section{Shannon's law}
\label{sec:orgf6e4163}
\begin{definition}
for every switching function \(f(x_1, \cdots, x_n)\) we get:\\
\begin{align*}
  f(x_1, \cdots, x_n) &= x_1 \cdot f(1, x_2, \cdots, x_n) + x_1' \cdots f(0, x_2, \cdots, x_n)\\
  f(x_1, \cdots, x_n) &= [x_1 + f(0, x_2, \cdots, x_n)] \cdot [x_1' + f(1, x_2, \cdots, x_n)]
\end{align*}
\begin{characteristic}
this law helps us move between \hyperref[sec:org03ef0dd]{DNF} and \hyperref[sec:org51d7cda]{CNF} forms without looking at the truth table of a function\\
\end{characteristic}
\begin{my_example}
\begin{align*}
  f(x, y, z) &= y'xz + y'z'x' + xz\\
             &= x(y'1z + y'z'1' + 1z) + x'(y'0z + y'z'0' + 0z)\\
             &= x(y'z + y'z'0 + z) + x'(0 + y'z' + 0)\\
             &= x(y'z + z) + x'(y'z')
\end{align*}
\end{my_example}
\end{definition}
\section{XOR}
\label{sec:orgf13d9a5}
\begin{tabular}{rrr}
\(x\) & \(y\) & \(x \oplus y\)\\
\hline
0 & 0 & 0\\
0 & 1 & 1\\
1 & 0 & 1\\
1 & 1 & 0\\
\end{tabular}
\begin{characteristic}
\begin{align}
  A \oplus B &= B \oplus A & \text{commutative}\\
  (A \oplus B) \oplus C &= A \oplus (B \oplus C) & \text{associative}\\
  A \bullet (B \oplus C) &= (AB) \oplus (AC) & \text{distributive}\\
  A \oplus \bar{A} &= 1,\ A \oplus 1 = \bar{A},\ A \oplus 0 = A\\
  \overline{A \oplus B} &= 1 \text{ only if } A = B
\end{align}
\end{characteristic}
\begin{characteristic}
if \(A \oplus B = C\) then:\\
\begin{align*}
  A \oplus C &= A \oplus (A \oplus B) = (A \oplus A) \oplus B = B\\
  B \oplus C &= B \oplus (A \oplus B) = (B \oplus B) \oplus A = A\\
  A \oplus B \oplus C &= A \oplus B \oplus (A \oplus B) = 0
\end{align*}
\end{characteristic}
\section{logic gate}
\label{sec:org182e608}
\subsection{AND gate}
\label{sec:orged54c3e}
\[
  t = a \bullet b
\]\\
\includesvg{/home/mahmooz/brain/out/bKKN1Bx}\\

\subsection{OR gate}
\label{sec:org3113977}
\[
  t = a + b + c
\]\\
\includesvg{/home/mahmooz/brain/out/JEh5DG2}\\

\subsection{NAND gate}
\label{sec:orgd71d708}
\[
  t = \overline{(a \bullet b \bullet c)}
\]\\
\includesvg{/home/mahmooz/brain/out/cINGfRg}\\

\subsection{NOR gate}
\label{sec:orgcd29f6d}
\[
  t = \overline{(a + b)}
\]\\
\includesvg{/home/mahmooz/brain/out/g8ASPWI}\\

\subsection{XOR gate}
\label{sec:org414bbfd}
\[
  t = a \oplus b
\]\\
\includesvg{/home/mahmooz/brain/out/FLVi83R}\\

\subsection{NOT gate}
\label{sec:org9d9f716}
\[
  b = \overline{a}
\]\\
\includesvg{/home/mahmooz/brain/out/f4X1VZI}\\

\section{logic circuit}
\label{sec:org1b137e5}
a combination of logical gates to represent a more complex switching function is called a logical circuit or a switching circuit\\

the variables of the function are called the \emph{inputs} of the circuit and the output of the function is the \emph{exit} of the circuit\\

the exit of a logical circuit depends only on the current combination of values of its input\\

\[
  A = T \bullet (H + W) + \overline{W} \bullet P
\]\\

\includesvg{/home/mahmooz/brain/out/PCw1P6a}\\

\subsection{Digital abstraction}
\label{sec:orgce1c66e}

how do we represent the values \(0\) and \(1\) using a digital logic circuits?\\

on one side the logic circuits have only 2 variables, voltage and time and they can have any value in a given range\\

and on the other side, boolean algebra has only 2 values which are the constants \(0\) and \(1\)\\

the \uline{solution}:\\

\begin{itemize}
\item voltage represents logical value, for example, a high voltage represents \(1\) and a low voltage represents \(0\), other voltages, for example a voltage between ``high'' and ``low'' are defined as not represnting of logical values\\
\item sometimes, it is agreed upon that the output of a circuit represents the value of the switching function, other times its agreed upon that it doesnt\\
\end{itemize}

\subsubsection{Logical levels}
\label{sec:orgd752773}

the two logical values \(0\) and \(1\) are represented in \emph{Electronic circuits} using logical levels.\\

an electronic circuit receives electrical energy from a power supply like a battery, the power supply is the high voltage \emph{Vplus} and the low voltage is called \emph{Vminus}\\

the levels of voltages denoted \(v\) in a circuit is limited in the range \(Vminus \leq v \leq Vplus\), the amplitude of voltages is defined as \(\Delta{v} = Vplus - Vminus\)\\

as an example:\\

\begin{align*}
  Vplus - 0.2\Delta{v} &= 4volt \leq v \leq 5volt = Vplus & \text{represents logical $1$}\\
  Vminus &= 0volt \leq v \leq 1volt = Vminus + 0.2\Delta{v} & \text{represents logical $0$}
\end{align*}
\end{document}